% ===== tr-08 · b14 · tr08a：第 5 章 p24--27（原文页 24--27）=====
% [原文 p.24]
记 $Y_r=T^n\times B(0,r)\subset T^n_{x_i}\times\R^n_{y_i}\simeq T^*T^n$，其中 $r\gg 1$ 固定。特殊拉格朗日纤维化最平凡的例子如下：其 CY 结构是平坦模型
\[
g=\sum (dx_i^2+dy_i^2),\qquad \omega=\sum dx_i\wedge dy_i,\qquad \Omega=\bigwedge (dx_j+\sqrt{-1}dy_j),
\]
而 Slag 纤维化就是到 $\R^n_{y_i}$ 因子的投影，即环面 $T^n\times\{y\}$ 都是特殊拉格朗日。Zhang 考虑一族 CY 结构 $(g_k,\omega_k,\Omega_k)$，它在 $Y_{2r}$ 上以 $C^\infty$ 意义收敛到 $(g,\omega,\Omega)$（这可由其有界截面曲率假设经由椭圆自举得到），且满足 $\omega_k\in[\omega]\in H^2(Y_{2r},\R)$。标准 $T^n$ 纤维的小形变可以表示为 $T^n$ 上的图：对 $y\in\R^n$ 以及 $T^n$ 上正交于调和 1-形式 $dx_1,\dots,dx_n$ 的 1-形式 $\sigma$，记
\[
L(y,\sigma)=\mathrm{Graph}(x\mapsto y+\sigma(x))\subset T^*T^n.
\]
$L(y,\sigma)$ 关于 $(g_k,\omega_k,\Omega_k)$ 是特殊拉格朗日的条件是
\begin{equation}
\omega_k|_{L(y,\sigma)}=0,\qquad \mathrm{Im}(e^{\sqrt{-1}\theta_k}\Omega_k)|_{L(y,\sigma)}=0, \tag{5.12}
\end{equation}
其中 $\theta_k$ 的选取使得 $\int_{T^n}e^{\sqrt{-1}\theta_k}\Omega_k>0$。Zhang 用扰动论证证明：对每个 $y\in B(0,\frac{3r}{2})$ 及 $k\ge k_0\gg 1$，存在唯一的 $\sigma=\sigma_{k,y}$，使得 $L(y,\sigma_{k,y})$ 解出式 (5.12) 且具有小范数界 $\|\sigma_{k,y}\|<\delta\ll 1$。他随后再用一次隐函数论证说明，这些特殊拉格朗日确实在 $Y_{3r/2}$ 的某个包含 $Y_r$ 的开子集上定义一个局部特殊拉格朗日 $T^n$-纤维化。

% [Tutor Note: p24 的 (12) 是本章第一个编号公式，按全章编号记为 (5.12)；原文写 "solves (12)"，指的就是本式，不要误以为另有一条未编号公式。式中的 $\Omega=\bigwedge(dx_j+\sqrt{-1}dy_j)$ 原文省略了外积下标 $j$，不要擅自补成 $\bigwedge_{j=1}^{n}$。]

\section{非阿基米德几何}
非阿基米德（NA）多复变势理论（pluripotential theory）是 Kähler 几何的密切类比。直观地说，
\begin{itemize}
\item NA 几何提供了描述复流形退化为实单纯/热带对象的自然语言；
\item 它系统地编码热带几何的组合学；
\item Kähler 几何中的通常概念，如函数、线丛、Kähler 度量、复 Monge--Ampère 测度，在 NA 几何中都有自然的（尽管看起来颇为奇特的）类比物；
\item Calabi 猜想的类比在 NA 情形下成立：可以解出 NA Monge--Ampère 方程；
\item 在附加假设下，NA Monge--Ampère 测度与实的 Monge--Ampère 测度一致。
\end{itemize}

% [Tutor Note: 这里的 "pluripotential theory" 译作「多复变势理论」（复数域上的势论），与第 3、4 章的 "plurisubharmonic/psh"（多重次调和）同源；不要译成「多潜能理论」。本节的 NA = non-archimedean（非阿基米德），后文一律沿用缩写 NA。]

% [原文 p.25]
我们将在下面说明，NA 多复变势理论的基本概念可以由如下启发式原则来理解：「NA 几何是复几何在混合拓扑下的极限」。作为一些不精确的直观，可以把非阿基米德几何近似地想象成 SYZ 底，而混合拓扑收敛大致相当于一个 SYZ 纤维化坍缩到其底。我们的视角深受 Boucksom 等人 [4][3][6][5] 的影响。

\subsection{Berkovich 空间、混合拓扑}
我们在第 5.3.2 节提到，对给定的极化代数退化，snc 模型的选取高度不唯一。关于如何提取不变量信息，有两种观点：
\begin{itemize}
\item 在双有理几何中，人们试图通过极小模型纲领寻找最优代表。通常这会离开 snc 世界，而需要因子化对数终端（divisorial log terminal）模型 [62][63]，但极小模型仍可能不唯一；
\item 在非阿基米德几何中，人们同时考察所有 snc 模型构成的塔，并取其对偶复形的形式极限，即所谓的 Berkovich 空间。
\end{itemize}
好的参考文献可见 [49, A]、[3, Appendix]、[6, chapter 2,3]。

Berkovich 的一个洞见是：把点看作乘法半范数，就得到一种类似于复流形的几何。设 $K\simeq\C((t))$ 配备其标准绝对值 $|\cdot|_0=e^{-\mathrm{ord}_t}$，其中 $\mathrm{ord}_t$ 是由消没阶（vanishing order）定义的值。其超度量性质
\[
|f+g|_0\le\max\{|f|_0,|g|_0\}
\]
赋予了这一主题「非阿基米德」之名。设 $X_K$ 是 $\mathrm{Spec}(K)$ 上的光滑、几何连通、射影概形；主要的例子来自把去心曲线（punctured curve）上的代数退化族 $X$ 作基变换。取 $X_K$ 的一个有限覆盖，其中各仿射开集形如 $U=\mathrm{Spec}(A)$，$A$ 是有限生成 $K$-代数。空间 $U^{an}$ 定义为所有延拓 $K$ 上绝对值的乘法半范数 $|\cdot|_x:A\to\R_{\ge 0}$ 构成的集合，并赋予使函数 $x\mapsto|f|_x$ 对任意 $f\in A$ 都连续的最弱拓扑。Berkovich 空间 $X_K^{an}$ 于是由粘合各个 $U^{an}$ 得到；这一记号代表「解析化」（analytification）。作为拓扑空间，$X_K^{an}$ 是紧的且 Hausdorff 的。在 CY 情形，$X_K^{an}$ 的点集描述意在编码关于 SYZ 纤维化底的信息；还存在一个自然的结构层，它编码关于复结构的信息 [49]。

% [Tutor Note: p25 的绝对值记号 $|\cdot|_0=e^{-\mathrm{ord}_t}$ 中，下标 $0$ 是取值（绝对值）的下标，$t$ 是 $\mathrm{ord}$ 的下标；文字层会抽成 "e−ordt"，务必按页图区分 $|\cdot|_0$ 与 $\mathrm{ord}_t$。另外 $K\simeq\C((t))$、$R\simeq\C[[t]]$ 用的是 $\simeq$（\simeq），不是 $\cong$。]

设 $R\simeq\C[[t]]$。$\mathrm{Spec}(R)$ 上模型的概念与代数曲线上的情形完全类似。$\mathrm{Spec}(R)$ 上 snc 模型的对偶相交复形 $\Delta_{\mathcal X}$ 可以通过两个自然映射与 $X_K^{an}$ 比较：

% [原文 p.26]
\begin{itemize}
\item 存在连续嵌入映射 $\mathrm{emb}:\Delta_{\mathcal X}\to X_K^{an}$。记 $\mathcal X_0=\sum b_iE_i$，每个除子 $E_i$ 通过消没阶 $\mathrm{ord}_{E_i}$ 定义 $\mathrm{val}_{E_i}=\frac{\mathrm{ord}_{E_i}}{b_i}$，于是 $e^{-\mathrm{val}_{E_i}}$ 是 $X_K^{an}$ 中的一点，称为除子型点（divisorial point）。更一般地，给定对应于 $E_J=\bigcap_{0}^{p}E_i$ 的面 $\Delta_J\subset\Delta_{\mathcal X}$ 内部的一点 $x=(x_0,\dots x_p)$，我们可以联系一个拟单项式赋值（quasi-monomial valuation）：把任意局部函数 $f$ 在 $E_J$ 附近展开成 Taylor 级数，
\[
f=\sum_{\alpha\in\N^{p+1}}f_\alpha z_0^{\alpha_0}\cdots z_p^{\alpha_p},\qquad f_\alpha\in K(E_J),
\]
则拟单项式赋值为
\[
\mathrm{val}_x(f)=\min\Big\{\sum_{0}^{p}\alpha_ix_i\Big|f_\alpha\ne 0\Big\}.
\]
于是 $x$ 给出一点 $e^{-\mathrm{val}_x}\in X_K^{an}$。我们将把 $\Delta_{\mathcal X}$ 视为 $X_K^{an}$ 的子集。特别地，本质骨架嵌入到 $X_K^{an}$ 中。
\item 存在连续收缩映射 $r_{\mathcal X}:X_K^{an}\to\Delta_{\mathcal X}$，它在 $\Delta_{\mathcal X}\subset X^{an}$ 上限制为恒等。任何点 $e^{-v}\in X_K^{an}$ 在 $\mathcal X$ 上都有中心。这就是唯一一个概形论的点 $\xi\in\mathcal X_0$，使得对 $f\in\mathcal O_{\mathcal X,\xi}$ 有 $|f|_x\le 1$，而对 $f\in m_{\mathcal X,\xi}$ 有 $|f|_x<1$。设 $J\subset I$ 是使 $\xi\in E_J$ 的极大子集。则 $r_{\mathcal X}(x)\in\Delta_{\mathcal X}$ 对应于一个拟单项式赋值，它在 $-\log|z_j|_x$（$j\in J$）上取相同的值。具体地说，应当把收缩映射理解为读出对数坐标。
\end{itemize}

换一个视角，若 $\mathcal X'$ 是 $\mathcal X$ 的爆破（blow up），则存在自然的单纯映射 $\Delta_{\mathcal X'}\to\Delta_{\mathcal X}$，它在 $\Delta_{\mathcal X}\subset\Delta_{\mathcal X'}$ 上是恒等。收缩映射 $X_K^{an}\to\Delta_{\mathcal X}$ 可以看作 $\mathcal X'$ 充分大时的形式极限。

% [Tutor Note: p26 的 $\xi$ 是「$\mathcal X_0$ 上的概形论点」而非一个 $\C$-点；$\mathcal O_{\mathcal X,\xi}$ 是局部环，$m_{\mathcal X,\xi}$ 是其极大理想（原文用斜体 $m$，不是 \mathfrak{m}）。两个条件 $|f|_x\le 1$ 与 $|f|_x<1$ 分别刻画「$x$ 在 $\xi$ 处取值不超过 1」与「$x$ 落在 $m_{\mathcal X,\xi}$ 对应的地方」，这正是中心（center）的定义。]

\begin{remark}{5.9}
收缩映射依赖于模型的选取。存在这样的例子：两个模型 $\mathcal X$ 与 $\mathcal X'$ 在 $X_K^{an}$ 中定义同一个子集 $\Delta_{\mathcal X}$，但收缩映射不同 [35, Appendix]。
\end{remark}

借助这些比较映射，Berkovich 空间 $X_K^{an}$ 同胚于 snc 模型对偶相交复形的逆向极限：
\[
X_K^{an}\simeq\varprojlim_{\text{snc 模型}}\Delta_{\mathcal X}.
\]
从概念上说，一个 snc 模型给出 Berkovich 空间的一个有限逼近。

% [Tutor Note: 逆向极限记号原文为 $\lim_{\longleftarrow}$（箭头在 "lim" 下方），等价于 \varprojlim；amssymb/amsmath 已载入，可直接用。不要把下标 "snc models" 写成限的上标。]

现在我们指出 NA 几何如何与复几何统一起来。考虑去心曲线上的代数退化 $X$。用 $|\cdot|$ 表示复数的通常绝对值。给定 $\C$-点 $z\in X_t$（$0<|t|\ll 1$），在 $X$ 的某个仿射坐标卡 $U=\mathrm{Spec}(A)$ 中，我们可以定义一个乘法半范数 $A\to\R_{\ge 0}$（注意并非非阿基米德！）
\begin{equation}
f\mapsto e^{-\log|f(z)|/\log|t|}=|f(z)|^{1/|\log|t||}. \tag{5.13}
\end{equation}

% [Tutor Note: (5.13) 右端指数是 $1/|\log|t||$（外层绝对值套住 $\log|t|$），因 $0<|t|\ll 1$ 时 $\log|t|<0$，故 $e^{-\log|f|/\log|t|}=|f|^{1/|\log|t||}$。文字层把它抽成 "1/∣log ∣t∣∣"，看起来像有两个绝对值竖线，实为嵌套：$|f(z)|^{\,1/|\log|t||}$。]

% [原文 p.27]
当一列点 $z$ 趋向 $t\to 0$ 时，对底上任意给定的亚纯函数 $f=\sum a_kt^k$，有
\[
\lim_{t\to 0}\log|f(z)|/\log|t|=\mathrm{ord}_0(f),
\]
这就是 $K$ 上标准的 NA 赋值。因此 $X_K^{an}$ 上的点是 $X_t$ 上由 $\C$-点定义的乘法半范数的自然极限。通过引入 $X\sqcup X_K^{an}$ 上的混合拓扑，可以把这一想法形式化，使得 $X_K^{an}$ 取代中心纤维的位置 [3, Appendix]。这时 $f\in A$ 诱导出 $X\sqcup X_K^{an}$ 上的局部连续函数。

第 5.3.1 节讨论的「混合」空间 $X\sqcup\Delta_{\mathcal X}$ 可以理解为一个有限逼近。给定一个 snc 模型 $\mathcal X$，取一列 $\C$-点 $q_k$ 趋向 $e^{-v}\in X^{an}$，它们在收缩映射 $r_{\mathcal X}$ 下的像为 $x=(x_0,\dots x_p)\in\Delta_J\subset\Delta_{\mathcal X}$。按定义，$q_k$ 集中在 $E_J$ 附近，而在局部坐标 $z_0,\dots z_p$ 中有 $\log|z_i(q_k)|/\log|t|\to v(z_i)=x_i$，这等价于 $\mathrm{Log}_{\mathcal X}(z_k)\to x=(x_0,\dots x_p)\in\Delta_J\subset\Delta_{\mathcal X}$。形式上，$X\sqcup X_K^{an}$ 上的拓扑就是取遍所有 snc 模型时 $X\sqcup\Delta_{\mathcal X}$ 的逆向极限。

% [Tutor Note: p27 中 $\mathrm{Log}_{\mathcal X}$ 是大写 Log（对数映射），与同页小写 $\log$ 的通常对数不同，二者不可混用；原文在此处切换记号是有意的。$\sqcup$ 为不交并（\sqcup），不是 $\cup$。]

\subsection{模型函数、度量、正性}
现在我们讨论 $X_K^{an}$ 上的函数、线丛与度量 [6]。给定 $\mathrm{Spec}(R)$ 上的模型 $\mathcal X$ 以及支撑在中心纤维 $\mathcal X_0$ 上的 Cartier 除子 $D$，我们可以通过设定
\[
\phi_D(x)=\max\{\log|f|_x : f\in\mathcal O_{\mathcal X}(D)\}
\]
来联系 $X_K^{an}$ 上的一个连续函数。对应 $D\mapsto\phi_D$ 按 $\Q$-线性性延拓。通过所有这样的模型与除子选取在 $\Q$-张成中得到的函数称为 $X_K^{an}$ 上的模型函数，它们构成 $C^0(X_K^{an})$ 的稠密子集。这类函数在对偶相交复形上的限制是分段仿射的。

% [Tutor Note: $\mathcal O_{\mathcal X}(D)$ 的花体 O 与花体 $\mathcal X$ 都不可省；文字层抽成 "OX (D)"，若照抄会失去「模型上的结构层」这一层含义。$\max\{\log|f|_x\}$ 的 $x$ 是 $|f|$ 的下标（点 $x$ 处的半范数取值），不是 $\log$ 的下标。]

为了理解其复几何意义，我们考虑由代数曲线 $S$ 上的 snc 模型作基变换得到的模型。除子 $D$ 在 snc 模型的全空间上规定了一类具有解析奇点的函数：
\[
\phi=\log|f|+C^\infty\ \text{函数},
\]
其中 $f$ 是 $D$ 的局部定义函数。当我们考虑对 $X_t$ 上限制作重新标度
\[
\phi_t=\frac{1}{\log|t|}\phi|_{X_t}
\]
时，在 $t\to 0$ 的极限中只有奇异项是相关的，且 $\phi_t$ 在混合拓扑下收敛到 $-\phi_D$。

我们通过 GAGA 原理来理解 $X_K^{an}$ 上的线丛：$X_K^{an}$ 上的线丛对应于概形 $X_K$ 上的线丛 $L$。$L$ 上的连续度量给每个局部截面 $s$ 指定 $X_K^{an}$ 的开子集上的一个非负连续局部函数 $\|s\|$，与层结构相容，使得 $\|fs\|(x)=|f|_x\|s\|(x)$，且当 $s$ 是 $L$ 的局部标架时 $\|s\|>0$。给定一个连续度量，$L$ 上任何其他连续度量都形如 $\|\cdot\|e^{-\phi}$（$\phi\in C^0(X^{an})$），这与 Hermite 度量与 Kähler 势之间的通常关系类似。

% [Tutor Note: 跨 p27/p28 的这一句 "any other continuous metric on $L$ is of the form $\|\cdot\|e^{-\phi}$ for some $\phi\in C^0(X^{an})$, analogous to..." 由本块（起于 p27）写完；p28 的 "As such $\phi$ is referred to as a potential function." 是下一句，归 tr08b 开头。另注意 $\|fs\|(x)=|f|_x\|s\|(x)$ 中 $f$ 是 $\C$ 值函数、$|f|_x$ 是点 $x$ 处的 NA 绝对值，不要写成模平方。]
